%% 50-usage.tex — Usage and Reusability

\section{Usage and Reusability}
\label{sec:usage}

\subsection{SPARQL Access}

A public read-only SPARQL 1.1 endpoint is available at \url{https://gemea.ise.fiz-karlsruhe.de/sparql}, backed by QLever.
Listings~1--3 illustrate representative query patterns; Listing~3 demonstrates cross-graph querying over \texttt{graph/ddbedm} and \texttt{graph/mocho}.

\begin{lstlisting}[language=SPARQL, basicstyle=\ttfamily\small, breaklines=true, caption={Objects associated with a given agent (GND URI).}]
PREFIX edm: <http://www.europeana.eu/schemas/edm/>
PREFIX dc:  <http://purl.org/dc/elements/1.1/>
SELECT ?cho ?title ?provider WHERE {
  ?cho dc:creator <http://d-nb.info/gnd/118540238> ; # Goethe, Johann Wolfgang von
       dc:title   ?title .
  ?agg edm:aggregatedCHO ?cho ;
       edm:dataProvider   ?provider .
  FILTER(isLiteral(?provider))
}
LIMIT 100
\end{lstlisting}

\begin{lstlisting}[language=SPARQL, basicstyle=\ttfamily\small, breaklines=true, caption={Objects associated with a given place.}]
PREFIX dc:      <http://purl.org/dc/elements/1.1/>
PREFIX dcterms: <http://purl.org/dc/terms/>
SELECT ?cho ?title ?date WHERE {
  ?cho dcterms:spatial <https://sws.geonames.org/2925533> ; # Frankfurt am Main
       dc:title         ?title ;
       dc:date          ?date .
  FILTER(STRSTARTS(?date, "180"))
}
LIMIT 100
\end{lstlisting}

\begin{lstlisting}[language=SPARQL, basicstyle=\ttfamily\small, breaklines=true, caption={Source type strings alongside aligned mocho classes.}]
PREFIX rdf: <http://www.w3.org/1999/02/22-rdf-syntax-ns#>
PREFIX dc:  <http://purl.org/dc/elements/1.1/>
PREFIX owl: <http://www.w3.org/2002/07/owl#>
SELECT ?mochoClass ?dctype (COUNT(DISTINCT ?cho) AS ?n) WHERE {
  GRAPH <https://gemea.ise.fiz-karlsruhe.de/graph/mocho> {
    ?mocho rdf:type ?mochoClass ;
           owl:sameAs ?cho .
  }
  GRAPH <https://gemea.ise.fiz-karlsruhe.de/graph/ddbedm> {
    ?cho dc:type ?dctype .
  }
}
GROUP BY ?mochoClass ?dctype ORDER BY DESC(?n) LIMIT 20
\end{lstlisting}


\subsection{Research Applications}

GEMEA's named graph structure, GND linking, provenance traces, and MCP interface support the following research use cases.

\paragraph{KG-grounded retrieval.}
The SPARQL endpoint and MCP interface support evaluation of KG-grounded RAG pipelines~\cite{pan2024unifying}, entity-linked question answering~\cite{lan2023complex}, and natural language question answering over structured data~\cite{usbeck2018qald}.

\paragraph{Entity embedding and link prediction.}
The entity graph (agents, places, timespans, providers, \texttt{mocho:Work} groupings, and domain- and modality-specific classes emitted by the mocho alignment) provides a large-scale benchmark for entity embedding and link prediction methods~\cite{wang2017kgesurvey}, including LLM-enhanced and contrastive approaches.
Named graph structure supports provenance-aware experiments that distinguish source triples from enrichment triples.

\paragraph{Explainability.}
Provenance-tracked named graphs and PROV-O\cite{lebo2013provo} traces on GND-linked triples support evaluation of KG-grounded explanation methods, including tracing a retrieval result back to its source institution, enrichment pathway, and match confidence.
The same structure supports evaluation of LLM-assisted enrichment provenance, where standard PROV-O traces are insufficient to capture model identity, confidence, and refinement type.


\paragraph{Metadata quality analysis.}
GEMEA exposes systematic encoding anomalies in source metadata that are invisible to record-level inspection but surface immediately as graph queries.
The following query identifies objects attributed to Goethe as \texttt{dc:creator} or \texttt{dc:contributor} in museum records, grouped by \texttt{dc:type} (see also Section~\ref{sec:quality}):
\begin{lstlisting}[language=SPARQL, basicstyle=\ttfamily\small, breaklines=true]
PREFIX rdf: <http://www.w3.org/1999/02/22-rdf-syntax-ns#>
PREFIX dc:  <http://purl.org/dc/elements/1.1/>
PREFIX edm: <http://www.europeana.eu/schemas/edm/>
SELECT ?dctype (COUNT(DISTINCT ?cho) AS ?n)
WHERE {
  GRAPH <https://gemea.ise.fiz-karlsruhe.de/graph/ddbedm> {
    { ?cho dc:creator    <http://d-nb.info/gnd/118540238> }
    UNION
    { ?cho dc:contributor <http://d-nb.info/gnd/118540238> }
    OPTIONAL { ?cho dc:type ?dctype }
    ?agg edm:aggregatedCHO ?cho ;
         edm:dataProvider  ?inst .
    ?inst rdf:type <http://ddb.vocnet.org/sparte/sparte006> .
  }
}
GROUP BY ?dctype ORDER BY DESC(?n)
\end{lstlisting}
The query returns 120 records from the museum sector (\texttt{sparte006}) across \texttt{dc:type} values including prints (59), stage productions (30), medals (5), and minting dies (2).
Within this set, 8 numismatic records (medals, minting dies, banknotes) show a clear misattribution: \texttt{dc:title} follows the pattern \emph{Maker: Johann Wolfgang von Goethe} (e.g., ``Bovy, Jean Fran\c{c}ois Antoine: Johann Wolfgang von Goethe''), where Goethe is the depicted subject but was recorded as \texttt{dc:creator} rather than the actual maker.
This encoding error, propagated from institutional source records through the DDB aggregator, is invisible without cross-institution graph queries.

\paragraph{Ontology evaluation.}
The \textbf{MOCHO}-aligned subgraph provides a grounded evaluation target for ontology alignment methods.
Because \texttt{graph/ddbedm} and \texttt{graph/mocho} encode the same objects under different schemas, the same query intent can be expressed against both graphs and compared on SPARQL execution time, result set size, and query complexity.
Typed properties and WEMI-level distinctions in the aligned graph should reduce query verbosity and improve selectivity relative to the flat EDM representation; divergence on these metrics gives a measurable alignment quality signal without a manually curated gold standard.

\subsection{Self-Hosting}
The full GEMEA stack is deployable via Docker Compose, with setup scripts for multiple configurations (full stack, endpoint-only, MCP proxy) provided in the repository.
Operators should expect at least 2\,TB of disk space for the QLever index and source data.
